% EXTRACTED FOR PAPER TWO (the hardware case study) -- PASS 48.1
% source     : paper/paper.tex
% blob       : 13590dbfe8119022ae95ff7a374910e30146ab8c
% lines      : 1554-1580 (1-indexed, inclusive)
% section    : Section 8
% prose words: 224
%
% Four Limitations items: the channel abstraction (which points at Section 6), the crossover not demonstrated on hardware, the computed-not-measured single-copy baseline, and the unverifiable Public Tier execution path.
%
% This is a VERBATIM COPY. The original is unmodified and still frozen.
% Regenerate with: experiments/pass48_extract_paper2.py

\textbf{The noise model is a channel abstraction.} Global depolarizing
and independent per-qubit channels are idealizations. The bias laws are
exact given the channel, and they say nothing about whether the channel
describes any particular device. Section 6 documents exactly where a
real device departs from them.

\textbf{The crossover was not demonstrated on hardware.} It sits at
\(n \approx 5\) to \(8\) under realistic noise, and
cross-session variation dominates at those sizes on the hardware available to
us. Section 6 is a test of the model, not a demonstration of the
advantage.

\textbf{The hardware single-copy baseline at \(n \geq 3\) is computed,
not measured}, for the economic reason given in Section 6.6. The one
point we could afford to measure, the \(n = 2\) anchor, missed its
prediction by 5.6 predicted standard deviations (Section 6.4). The computed points are
therefore model output on a device that our own anchor shows the model
does not currently track, and nothing in Section 6 should be read as a
hardware measurement of the single-copy route.

\textbf{The execution path on the Public Tier is not independently
verifiable by us.} Jobs were routed through a decentralized network with
validator spot-checking rather than direct vendor access (Section 6.1).
We cannot exclude this as a contributor to the session-to-session
variability of Section 6.4. It does not affect the readout and gate
characterizations, which are internally consistent, but a replication
with direct vendor access would be worth running.
